\section{Research Hypothesis}


This study investigates whether prompt engineering can change the traditional
algorithm learning and problem-solving workflow used by computer science
students.


The research focuses on the following hypotheses.



\subsection{H1: Prompt Optimization Improves Solution Correctness}


\textbf{Hypothesis H1:}

Structured prompt engineering improves the correctness of LLM-generated
algorithm solutions compared with direct prompting.


The hypothesis is evaluated by comparing the pass rate of generated solutions
under baseline and optimized prompt strategies.



\subsection{H2: Prompt Engineering Improves Algorithm Selection}


\textbf{Hypothesis H2:}

Optimized prompts improve the ability of LLMs to select appropriate data
structures and algorithms.


This hypothesis evaluates whether additional prompting guidance leads to:

\begin{itemize}

\item Better algorithm pattern selection.

\item More suitable data structures.

\item Improved computational complexity.

\end{itemize}



\subsection{H3: AI-Assisted Workflow Reduces Traditional Algorithm Practice Dependency}


\textbf{Hypothesis H3:}

AI-assisted programming workflows may reduce the amount of manual algorithm
memorization required for software engineering problem solving.


Instead of relying only on extensive LeetCode practice, future engineers may
benefit from developing skills in:

\begin{itemize}

\item Prompt engineering.

\item AI collaboration.

\item Solution verification.

\item System-level thinking.

\end{itemize}



\subsection{Research Questions}


Based on these hypotheses, the study investigates three research questions:


\textbf{RQ1:}

Can optimized prompts improve LLM algorithm-solving accuracy?



\textbf{RQ2:}

Can prompt engineering improve algorithm and data structure selection?



\textbf{RQ3:}

How might AI-assisted workflows change the skills required for future software
engineers?


